\metatitle{Polynomial recurrences}
\metadescription{Spaces of polynomial recurrences, exact finite verification,
first-order derivative criteria, repeated poles, limiting zeros, and the trade-off between
recurrence order and derivative order.}
\metakeywords{polynomial recurrence, derivative recurrence, Ore algebra,
rational generating function, recurrence space, repeated poles, Beraha--Kahane--Weiss}

\section[polynomialRecurrences]{Polynomial recurrences}

A polynomial sequence can have several useful recurrences. A relation using
many consecutive polynomials may be replaced by one using fewer polynomials
and their derivatives. We study when this is possible, how to describe all
relations with prescribed coefficient bounds, and how to prove an identity
found by exact linear algebra.

Throughout, $P_n(x)\in\setQ[x]$ and $n\geq0$. We write
\[
 \mathbf{P}=(P_0,P_1,\dotsc),\qquad
 F(x,t)\coloneqq\sum_{n\geq0}P_n(x)t^n.
\]
A linear \defin{differential--difference equation} of order $r$ in the index
and derivative order at most $m$ has the form
\[
 D(n,x)P_{n+r}(x)
 =\sum_{j=0}^{r-1}\sum_{s=0}^{m}A_{j,s}(n,x)P_{n+j}^{(s)}(x),
 \qquad D\ne0.
\]
The superscript denotes differentiation in $x$. For $m=0$ this is an ordinary
linear recurrence. An identity determines the next polynomial
only at indices where $D(n,x)$ is not the zero polynomial in $x$.
We shall distinguish polynomial coefficients with prescribed degree bounds,
arbitrary rational coefficients, and the restriction $D=D(n)$.

The general operator framework is the $\partial$-finite theory of
\name{Frédéric Chyzak} and \name{Bruno Salvy}
\cite[Sec.~2]{ChyzakSalvy1998}; see also \name{Christoph Koutschan}'s
account of the holonomic systems approach \cite{Koutschan2009Thesis}.
The criteria below are direct consequences of linear algebra and rational
generating functions. Real-rootedness is an
\hyperref[derivativeRecurrenceRealRoots]{additional application}.

\subsection[recurrenceTradeoffExample]{Reducing recurrence order using derivatives}

Consider
\[
 P_n(x)=1+x^n+(x+1)^n.
\]
It satisfies three different recurrences:
\begin{align*}
 P_{n+3}
  &=(2x+2)P_{n+2}-(x^2+3x+1)P_{n+1}+x(x+1)P_n,\\
 2nP_{n+2}
  &=3n(x+1)P_{n+1}-n(3x+1)P_n-x(x-1)(x+1)P_n',\\
 n(n-1)P_{n+1}
  &=n(n-1)P_n+(n-1)x(3x+1)P_n'-2x^2(x+1)P_n''.
\end{align*}
Each identity holds for $n\geq0$, as we see by substituting the three
summands separately.

\begin{tabular}{ccc}
 Recurrence order & Derivative order & Leading coefficient\\
 3 & 0 & $1$\\
 2 & 1 & $2n$\\
 1 & 2 & $n(n-1)$
\end{tabular}

\section[boundedRecurrenceSpaces]{Recurrences with bounded coefficient degrees}

Suppose the sequence is specified by $P_0,\dotsc,P_{k-1}$ and
\[
 P_{n+k}=\sum_{i=0}^{k-1}c_i(x)P_{n+i},\qquad c_i(x)\in\setQ[x].
\]
Such recurrences arise, for example, from the
\hyperref[transferMatrixMethod]{transfer matrix method}.

Let $S$ shift the index, so $(S\mathbf{P})_n=P_{n+1}$, and let
$\partial_x$ differentiate each polynomial. Fix a nonempty finite set of allowed
monomials and consider operators
\[
 L=\sum_{j,s,u,v}\gamma_{j,s,u,v}x^v n^u S^j\partial_x^s,
 \qquad \gamma_{j,s,u,v}\in\setQ.
\]
All indices in this sum are nonnegative. Coefficients act on the left, so
$(nS\mathbf{P})_n=nP_{n+1}$. These operators form a finite-dimensional vector space.

\begin{theorem}[Finite verification]
\label{recurrenceFiniteVerification}
Put $h=1+\max(u+s)$ over the allowed monomials. Then
\[
 L\mathbf{P}=\mathbf{0}
 \quad\Longleftrightarrow\quad
 (L\mathbf{P})_n=0\text{ for }0\leq n<hk.
\]
Consequently, all recurrence operators in the prescribed space form the
kernel of a finite matrix over $\setQ$.
\end{theorem}

\begin{proof*}
Write
\[
 F(x,t)=\frac{N(x,t)}{Q(x,t)},
 \qquad Q(x,0)=1,\quad \deg_tN<k,\quad\deg_tQ\leq k.
\]
Every shifted sequence has a generating function
$F_j(x,t)\coloneqq\sum_{n\geq0}P_{n+j}(x)t^n=N_j/Q$ with
$\deg_tN_j<k$. Multiplication by $n$ corresponds to
$\theta=t\partial_t$. Thus the residual generating function is
\[
 \sum_{n\geq0}(L\mathbf{P})_nt^n
 =\sum_{j,s,u,v}\gamma_{j,s,u,v}x^v\theta^u\partial_x^sF_j.
\]
Putting this expression over $Q^h$ gives a numerator of $t$-degree less than
$hk$: each differentiation or application of $\theta$ raises the needed
power of $Q$ by at most one and the numerator degree by at most $k$.
Since $Q(x,0)=1$, its first $hk$ coefficients vanish if and only if that
numerator is zero.
\end{proof*}

To compute the matrix, substitute an unknown coefficient for each allowed
monomial, evaluate these $hk$ residuals, and equate coefficients of $x$.
For $M$ unknowns and $R$ scalar equations, exact elimination uses
$\mathcal{O}(RM^2+M^3)$ rational field operations; coefficient sizes affect bit complexity.

A recurrence with $D\ne0$ exists in the prescribed space if and only if
the kernel has nonzero projection to the $D$ coordinates. Vectors with
$D=0$ are relations among the earlier polynomials and their derivatives.

\subsection[recurrencePropagation]{A recurrence for the residual}

There is also a finite induction proof. Put
\[
 \chi(z)\coloneqq z^k-\sum_{i=0}^{k-1}c_i(x)z^i,\qquad T\coloneqq\chi(S),
\]
and let $\mathsf N$ multiply a sequence by its index. The commutators satisfy
\[
 [T,\mathsf N]=(z\chi_z)(S),\qquad
 [\partial_x,T]=\chi_x(S).
\]
Both right-hand sides commute with $T$. For a sequence $\mathbf{f}$ with
$T^q\mathbf{f}=\mathbf{0}$, the commutator identities imply
$T^{q+1}(\mathsf N\mathbf{f})=\mathbf{0}$ and
$T^{q+1}(\partial_x\mathbf{f})=\mathbf{0}$.
Shifts and multiplication by $x$ commute with $T$, so every allowed residual
satisfies
\[
 \chi(S)^h(L\mathbf{P})=\mathbf{0}.
\]
This is a monic recurrence of order $hk$. Its first $hk$ zero polynomials
therefore force all subsequent residuals to vanish. The coefficients of
$\chi(z)^h$ and the initial residuals determine a proof by recurrence induction.

\section[recurrenceGeneratingFunctions]{First derivatives and generating functions}

For a proposed identity
\[
 (d_0+d_1n)P_{n+1}
 =(A_0(x)+nA_1(x))P_n+B(x)P_n',
\]
where $d_0,d_1\in\setQ$ and $A_0,A_1,B\in\setQ[x]$, summation gives the
following equivalent equation:
\[
 t(d_1-tA_1)F_t-tBF_x+(d_0-d_1-tA_0)F=(d_0-d_1)P_0.
\]
This is a first-order linear partial differential equation. If the derivative
coefficient is instead $B_0(x)+nB_1(x)$, the equation contains the additional
term $-t^2B_1F_{xt}$ and is generally of second order.

There is an exact criterion involving the reduced denominator. Define
\[
 V=t(d_1-tA_1)\partial_t-tB\partial_x,\qquad
 C=d_0-d_1-tA_0,\qquad H=(d_0-d_1)P_0.
\]

\begin{proposition}[Denominator and numerator criterion]
Suppose $F=N/Q\ne0$ with $N,Q\in\setQ[x,t]$ coprime. The proposed
first-derivative identity holds if and only if
\[
 Q\mid V(Q),\qquad
 V(N)+\left(C-\frac{V(Q)}{Q}\right)N=HQ.
\]
\end{proposition}

\begin{proof*}
After multiplying the generating-function equation by $Q^2$, it becomes
\[
 Q\,V(N)-N\,V(Q)+CNQ-HQ^2=0.
\]
Reduction modulo $Q$ and coprimality give $Q\mid V(Q)$.
Division by $Q$ gives the second condition, and the steps are reversible.
\end{proof*}

The pole curves are invariant under the vector field $V$, and the numerator
condition records the initial polynomials.

\section[recurrenceSpectralModel]{Rational coefficients and repeated poles}

We now allow coefficients in $\setQ(n,x)$ and consider identities for all
sufficiently large indices. Divide the reduced rational generating function
into a polynomial part $U(x,t)$ and a proper rational part. Partial fractions
over a finite splitting field $E$ of the characteristic polynomial over
$\setQ(x)$ give
\[
 P_n^{\mathrm{tail}}=\sum_{i=1}^d a_i(n,x)\lambda_i(x)^n.
\]
Here $P_n^{\mathrm{tail}}$ is the coefficient of $t^n$ in the proper rational
part, so $P_n^{\mathrm{tail}}=P_n$ for all sufficiently large $n$.
The $\lambda_i$ are distinct nonzero characteristic roots, and
$a_i\in E[n]$ is nonzero. A pole of multiplicity $m_i$ gives
$\deg_n a_i=m_i-1$, since
\[
 [t^n](1-\lambda_i t)^{-m_i}
 =\binom{n+m_i-1}{m_i-1}\lambda_i^n.
\]
On $E(n)^d$, with amplitude vector $\mathbf{a}=(a_i)$, define operators
on $\mathbf{v}=(v_i)$ by
\[
 (\mathsf S\mathbf{v})_i=\lambda_i v_i(n+1,x),\qquad
 (\nabla\mathbf{v})_i=\partial_xv_i+n\frac{\lambda_i'}{\lambda_i}v_i.
\]
Differentiation includes the algebraic dependence on $x$ in $E$.
These operators commute and represent shift and differentiation of the
corresponding exponential sums.

\begin{theorem}[Rank criterion]
\label{recurrenceRankCriterion}
Let $1\leq r\leq d$ and $m\geq0$. There is an eventual identity
\[
 P_{n+r}=\sum_{j=0}^{r-1}\sum_{s=0}^{m}A_{j,s}(n,x)P_{n+j}^{(s)},
 \qquad A_{j,s}\in\setQ(n,x),
\]
if and only if
\[
 \mathsf S^r\mathbf{a}\in
 W_{r,m}\coloneqq\operatorname{span}_{E(n)}
 \{\nabla^s\mathsf S^j\mathbf{a}:0\leq j<r,\ 0\leq s\leq m\}.
\]
If such an identity exists at any derivative order, one exists with
$m\leq d-r$.
\end{theorem}

\begin{proof*}
Exponential-polynomial sequences with distinct nonzero bases are independent,
so an identity is equivalent to equality of its component vectors.
A solution over the splitting field descends to $\setQ(n,x)$ by averaging
its Galois conjugates.

The vectors
$\mathbf{a},\mathsf S\mathbf{a},\dotsc,\mathsf S^{d-1}\mathbf{a}$ are independent.
Indeed, the leading coefficient in $n$ of their coordinate determinant is
the product of the leading coefficients of the $a_i$ times the nonzero
Vandermonde determinant in the $\lambda_i$. Hence $\dim W_{r,0}=r$.
If $W_{r,m+1}=W_{r,m}$, this space is closed under $\nabla$, including
derivatives of its scalar coefficients, and all later spaces agree.
Otherwise the dimension increases. There can be at most $d-r$ increases.
\end{proof*}

For identities holding at all sufficiently large indices, the least order
of a recurrence without derivatives is $d$ with coefficients in $\setQ(n,x)$,
and $\sum_i m_i$ with coefficients in $\setQ(x)$.

For example, $(n+1)x^n$ has a double pole at $t=1/x$ and satisfies both
\[
 P_{n+2}=2xP_{n+1}-x^2P_n,
 \qquad (n+1)P_{n+1}=(n+2)xP_n.
\]

\subsection[recurrenceAnnihilatingIdeal]{The annihilating ideal}

In the \defin{Ore algebra} $\setQ(n,x)\langle S,\partial_x\rangle$,
operators satisfy
\[
 S f(n,x)=f(n+1,x)S,\qquad
 \partial_x f=f\partial_x+\partial_x(f),\qquad S\partial_x=\partial_xS.
\]
The operators annihilating the sequence eventually form its
\defin{annihilating ideal}. Expressing $\mathsf S^d\mathbf{a}$ and $\nabla\mathbf{a}$ in
the shift basis gives two generators:
\[
 S^d-\sum_{j=0}^{d-1}r_j(n,x)S^j,\qquad
 \partial_x-\sum_{j=0}^{d-1}b_j(n,x)S^j.
\]
To see this, reduce derivatives using the second relation and then shifts
using the first. The remaining operator is a linear combination of
$1,S,\dotsc,S^{d-1}$, and independence forces an annihilating remainder
to be zero.

The polynomial part $U$ records the finite initial contribution discarded
by eventual equality. After clearing denominators, a polynomial-coefficient
tail identity holds from $n=0$ exactly when it also annihilates the finite
coefficient sequence of $U$. Checking $0\leq n\leq\deg_tU$ suffices,
since forward shifts and $x$-derivatives cannot enlarge this support.

\section[recurrenceLimitingZeros]{Beraha--Kahane--Weiss and limiting zeros}

The same characteristic roots also control where zeros accumulate as
$n\to\infty$. This is the content of the Beraha--Kahane--Weiss theorem
\cite{Beraha75limits}, also used in the
\hyperref[berahaKahaneWeiss]{transfer-matrix setting}.
We now let $x$ be complex and work on a domain where the characteristic roots
and amplitudes are single-valued analytic functions.

\begin{theorem}[Beraha--Kahane--Weiss, analytic form]
\label{recurrenceBkwTheorem}
Let $d\geq2$ and let $a_i,\lambda_i$ be analytic on a connected open set
$\Omega\subseteq\setC$, with none identically zero. Suppose
$\lambda_i\not\equiv\omega\lambda_j$ for $i\ne j$ and every constant
$|\omega|=1$. Set
\[
 P_n(x)=\sum_{i=1}^d a_i(x)\lambda_i(x)^n.
\]
Call $i$ \defin{dominant} at $x$ if
$|\lambda_i(x)|=\max_j|\lambda_j(x)|$.
A point $x\in\Omega$ is a limit of zeros $x_k$ of $P_{n_k}$ with
$n_k\to\infty$ if and only if either
\begin{itemize}
\item exactly one index $i$ is dominant and $a_i(x)=0$; or
\item at least two indices are dominant.
\end{itemize}
Every neighborhood of such a point contains a zero of $P_n$ for all
sufficiently large $n$.
\end{theorem}

In \cite[Thm.~1.5]{Sokal2004ChromaticRoots}, the condition on constant
unit-modulus ratios is imposed on branches that are dominant on a set
with nonempty interior.

\begin{example}
For $P_n(x)=1+x^n+(x+1)^n$, all amplitudes equal $1$, so there are no
isolated limits arising from amplitude zeros. The limiting set is
\begin{align*}
 &\{x:|x|=1,\ |x+1|\leq1\}\\
 {}\cup{}&\{x:|x+1|=1,\ |x|\leq1\}\\
 {}\cup{}&\{x:|x|=|x+1|\geq1\}.
\end{align*}
The first two sets are circular arcs. The third consists of the two rays
\[
 \operatorname{Re}x=-\tfrac12,\qquad
 |\operatorname{Im}x|\geq\tfrac{\sqrt3}{2}.
\]
All three branches meet at $(-1\pm i\sqrt3)/2$. These equations follow by
comparing the moduli of $1,x,x+1$. In particular, the nonreal limiting points
force nonreal zeros for all sufficiently large $n$.
\end{example}

The derivative-recurrence criterion uses the values and logarithmic
derivatives of the characteristic roots; BKW uses their relative moduli.
Both depend on the amplitudes, which are determined by the initial polynomials.

For repeated poles the amplitudes $a_i(n,x)$ depend polynomially on $n$,
so the displayed theorem does not apply directly. A common scalar factor
in $n$ can be removed for large $n$; general polynomial-in-$n$ amplitudes
require an extension.

\section[firstDerivativeCriterion]{When does one derivative suffice?}

We return to the amplitudes $a_i(n,x)$ from
\hyperref[recurrenceSpectralModel]{the repeated-pole expansion}. Put
\[
 \eta_i=\frac{\partial_xa_i}{a_i}
       +n\frac{\lambda_i'}{\lambda_i},\qquad
 \rho_i=\lambda_i\frac{a_i(n+1,x)}{a_i(n,x)}.
\]

\begin{proposition}[First-order criterion, including repeated poles]
Suppose $d\geq2$. An identity $P_{n+1}=AP_n+BP_n'$ with
$A,B\in\setQ(n,x)$ exists eventually if and only if the $\eta_i$ are
not all equal and there are $A,B$ such that
\[
 \rho_i=A+B\eta_i\qquad\text{for every }i.
\]
Equivalently,
\[
 \operatorname{rank}[\mathbf{a},\nabla\mathbf{a},\mathsf S\mathbf{a}]
 =\operatorname{rank}[\mathbf{a},\nabla\mathbf{a}]=2.
\]
When this identity exists, its coefficients $A,B$ are unique.
For any $i,j$ with $\eta_i\ne\eta_j$,
\[
 B=\frac{\rho_i-\rho_j}{\eta_i-\eta_j},\qquad
 A=\rho_i-B\eta_i.
\]
\end{proposition}

\begin{proof*}
Divide each component equation by $a_i$.
The vectors $\mathbf{a}$ and $\mathsf S\mathbf{a}$ are independent by the
shift-basis argument, so a solution requires independent vectors
$\mathbf{a},\nabla\mathbf{a}$.
The two displayed scalar equations then determine $A,B$ uniquely.
\end{proof*}

For simple poles the amplitudes are independent of $n$, and the criterion
has a more explicit form. Write $u_i=a_i'/a_i$ and
$v_i=\lambda_i'/\lambda_i$. A first-order differential recurrence exists
if and only if
\[
 u_i=\alpha_0+\beta_0\lambda_i,\qquad
 v_i=\alpha_1+\beta_1\lambda_i
\]
for all $i$, with $\alpha_0,\alpha_1,\beta_0,\beta_1\in\setQ(x)$ and
$(\beta_0,\beta_1)\ne(0,0)$. In that case
\[
 A=-\frac{\alpha_0+n\alpha_1}{\beta_0+n\beta_1},\qquad
 B=\frac{1}{\beta_0+n\beta_1}.
\]
Indeed, $\lambda_i=A+B(u_i+nv_i)$ is equivalent to these conditions after
comparing coefficients of $n$ modulo the span of the vectors
$\mathbf{1}=(1,\dotsc,1)$ and $\boldsymbol{\lambda}=(\lambda_i)$.
In particular, all characteristic roots satisfy the same Riccati equation
$\lambda_i'=\alpha_1\lambda_i+\beta_1\lambda_i^2$.

Repeated poles still give a degree bound. If
$\mu\coloneqq\max_i(m_i-1)$ and a first-order differential recurrence exists,
it has a polynomial form
\[
 D(n,x)P_{n+1}=\widetilde{A}(n,x)P_n+\widetilde{B}(n,x)P_n'
\]
with index degrees at most $2\mu+1,2\mu+1,2\mu$, respectively, for the proper tail.
To obtain it, put $p(z)=\prod_i(z-\lambda_i)$ and use a basis of
$\setQ(x)[z]/(p(z))$. The vectors
$\mathbf{a},\nabla\mathbf{a},\mathsf S\mathbf{a}$ have index degrees at
most $\mu,\mu+1,\mu$. A nonzero two-by-two minor and Cramer's rule give the
stated bounds; clearing denominators in $x$ preserves them.

\section[recurrenceCoefficientRestrictions]{Why the coefficient class matters}

More derivatives do not always shorten a recurrence. For
\[
 P_n(x)=(1+2^n)x^n,
\]
every derivative is a rational multiple of $P_n$, whereas $P_{n+1}$ is not.
The two characteristic roots $x$ and $2x$ have the same logarithmic derivative.
Thus no finite number of derivatives of $P_n$ can express $P_{n+1}$
with rational coefficients.

Requiring the coefficient of $P_{n+1}$ to depend only on $n$ creates another
obstruction. For each integer $\ell\geq1$, the sequence
\[
 P_n(x)=1+(2+x^\ell)^n
\]
has two simple poles. Over $\setQ(n,x)$, one derivative suffices:
\[
 P_{n+1}=P_n+
 \frac{(1+x^\ell)(2+x^\ell)}{n\ell x^{\ell-1}}P_n'.
\]
With polynomial right-hand coefficients and a nonzero leading coefficient
$D(n)$, the least derivative order is instead exactly $\ell$.

\begin{proof*}
Put $\lambda=2+x^\ell$ and $d_s=\lambda^{-n}\partial_x^s\lambda^n$,
viewed in $\setQ(n,x)$. The constant exponential component forces the
coefficient of $P_n$ to equal $D(n)$. The other component requires
\[
 D(n)(\lambda-1)=\sum_{s=1}^m C_s(n,x)d_s.
\]
If $m<\ell$, every $d_s$ vanishes at $x=0$, while $\lambda-1=1$ there.

For the upper bound, the polynomials $\lambda^\ell d_1$ and
$\lambda^\ell d_\ell$ are coprime in $\setQ(n)[x]$.
The roots of the first can only be $0$ or roots of $\lambda$.
At $0$ the second equals $\ell!n2^{\ell-1}$; at a root $\alpha$ of
$\lambda$ it equals
$n(n-1)\dotsm(n-\ell+1)\lambda'(\alpha)^\ell$.
Both are nonzero in $\setQ(n)$ or its algebraic extension.
Bézout's identity therefore expresses $\lambda-1$ as a polynomial-in-$x$
combination of $d_1,d_\ell$. Clearing denominators in $n$ proves existence.
\end{proof*}

For prescribed coefficient-degree bounds, we can compute these recurrences
using the matrix in the \hyperref[recurrenceFiniteVerification]{finite verification theorem}.

\section[recurrenceComputation]{Computing recurrences}

% Related Rust: polytool/src/recurrence/space.rs at source commit 6d10a1a.
The \href{https://poly.symmetricfunctions.com/#polynomial-sequence}{Polytool
recurrence finder} provides comparison and complete-space modes.
A useful report records the allowed coefficient degrees and derivatives,
the starting index, a basis of the recurrence space, and formulas with
$D\ne0$.

With a given recurrence and initial polynomials, the
\hyperref[recurrenceFiniteVerification]{finite verification theorem}
proves the identities and completeness within the specified bounds.
With finitely many input polynomials alone, a successful fit is an identity
on those inputs. If every kernel vector has $D=0$, no recurrence
with $D\ne0$ and the allowed monomials can hold at every index: it
would have to satisfy the same finite constraints. A computation stopped by
a resource bound gives no such nonexistence conclusion.

For given recurrences whose coefficients depend polynomially on $n$, the
rational generating-function argument need not apply. One can instead differentiate
the given recurrence and form a finite state of shifted polynomials and
their derivatives. A scalar recurrence for the residual, together with
checks at singular indices and enough initial zeros, supplies a finite
verification argument. Its bound depends on that state system.

The \hyperref[polynomialRootIntro]{real-rootedness pages} discuss criteria
for proving real-rootedness and interlacing, including the
\hyperref[maWangDerivativeCriterion]{Ma--Wang derivative criterion}.

\begin{example}
\label{derivativeRecurrenceRealRoots}
Let $a<b$ and suppose
\[
 D(n)P_{n+1}=A(n,x)P_n+\kappa_n(x-a)(x-b)P_n',
\]
where $D(n)>0$, $\kappa_n\geq0$ and $A(n,x)$ is a real polynomial in $x$.
Suppose $P_n$ has all its roots in $[a,b]$, both $P_n,P_{n+1}$ have positive
leading coefficients, and $\deg P_{n+1}\in\{\deg P_n,\deg P_n+1\}$.
Then $P_{n+1}$ is real-rooted and $P_n\interl P_{n+1}$ by the
\hyperref[maWangDerivativeCriterion]{Ma--Wang criterion}: at every root
$\xi$ of $P_n$, we have $\kappa_n(\xi-a)(\xi-b)/D(n)\leq0$.
\end{example}
